\section{Summary of Analytical Solutions} \label{cha:Analytical_Summary}
%
% A hyperbolic choice of q(x)
\begin{case} \label{cas:q_1st_Degree_Denom}
	Consider now the simple \gls{bar} problem with a hyperbolic $q \left( x \right)$ and a polynomial load function i.e.,
	\begin{subequations} \label{eqn:q_1st_Degree_Denom_Problem}
		\begin{align}	\label{eqn:q_1st_Degree_Denom_Diff_Eq}
			-\frac{ d }{ dx } \Biggl[ \gls{q} \glsentryname{uPrime} \Biggr] &= \gls{fDegreen}, \quad x \in \left( 0, 1 \right),	\\
			\gls{q} &= \frac{ e_1 }{ x + e_2 },  \label{eqn:q_1st_Degree_Denom_q_Func} \\
			f_{n} \left( x \right) &= \sum_{i = 0}^{n}c_{i}x^{i}, \label{eqn:q_1st_Degree_Denom_Load_Func}
		\end{align}
		with boundary conditions
		\begin{align}
			\gls{u} \Big \vert_{x = 0}  &= \gls{delta},  \label{eqn:q_1st_Degree_Denom_Left_Boundary} \\
			\Biggl[ \gls{q} \glsentryname{uPrime} \Biggr]_{x = 1} &= P. \label{eqn:q_1st_Degree_Denom_Right_Boundary}
		\end{align}
	\end{subequations}
	For $q$ to remain finite and strictly positive on $[0,1]$, the admissible parameters satisfy either $e_1>0$ and $e_2>0$, or $e_1<0$ and $e_2<-1$.
	\Gls{analSol} of \cref{cas:q_1st_Degree_Denom} for a general second degree polynomial load is
	\begin{subequations}	\label{eqn:q_1st_Degree_Denom_Solution}
		\begin{equation}
			u \left( x \right) = \sum_{i = 0}^{5} d_i x^i,
		\end{equation}
		where
		{\allowdisplaybreaks
		\begin{align}
			d_0 &= \gls{delta},	\\
			d_1 &= \frac{e_2}{e_1} P + \frac{e_2}{e_1} c_0 + \frac{e_2}{2 e_1} c_1 + \frac{e_2}{3 e_1} c_2,	\\
			d_2 &= \frac{1}{2 e_1} P + \frac{1 - e_2}{2 e_1} c_0 + \frac{1}{4 e_1} c_1 + \frac{1}{6 e_1} c_2,	\\
			d_3 &= -\frac{1}{3 e_1} c_0 - \frac{e_2}{6 e_1} c_1,	\\
			d_4 &= -\frac{1}{8 e_1} c_1 - \frac{e_2}{12 e_1} c_2, \\
			d_5 &= -\frac{1}{15 e_1} c_2.
		\end{align}}	% end of \allowdisplaybreaks
	\end{subequations}
\end{case}
%
\begin{remark}
	Note that by setting either $c_1 = c_2 = 0$ or $c_2 = 0$ in \cref{eqn:q_1st_Degree_Denom_Solution} we automatically get
	the solutions of linear and uniform load, respectively. The solution in \cref{eqn:Classical_Bar} is thus a special
	case of \cref{eqn:q_1st_Degree_Denom_Solution}.
\end{remark}
%
% Solution of the classical problem with q(x) = e_1 / (x^2 + e_2)
\begin{case}	\label{cas:q_2nd_Degree_Denom}
	Consider now \cref{prb:Normalized_Problem}, with another choice of $q \left( x \right)$, i.e.,
	\begin{subequations} \label{eqn:q_2nd_Degree_Denom_Problem}
		\begin{align}
			-\frac{d}{dx} \biggl[ q \left( x \right) \frac{du}{dx} \biggr] &= f_{n} \left( x \right),
				\quad x \in \left( 0, 1 \right), \label{eqn:q_2nd_Degree_Denom_Diff_Eq}	\\
			\gls{q} &= \frac{ e_1 }{x^{2} + e_2 },  \label{eqn:q_2nd_Degree_Denom_q_Func} \\
			\gls{fDegreen} &= \sum_{i = 0}^{n} c_i x^i, \label{eqn:q_2nd_Degree_Denom_Load_Func}
		\end{align}
		with boundary conditions
		\begin{align}
			u \left( x \right) \Big \vert_{x = 0} &= \gls{delta},  \label{eqn:q_2nd_Degree_Denom_Left_Boundary} \\
			\biggl[ q \left( x \right) \frac{du}{dx} \biggr]_{x = 1} &= P. \label{eqn:q_2nd_Degree_Denom_Right_Boundary}
		\end{align}
	\end{subequations}
	As in the linear-denominator case, strict positivity and finiteness on $[0,1]$ require either $e_1>0$ and $e_2>0$, or $e_1<0$ and $e_2<-1$.
	\Gls{analSol} of \cref{cas:q_2nd_Degree_Denom} for a polynomial load up to second degree is
	\begin{subequations}	\label{eqn:q_2nd_Degree_Denom_Solution}
		\begin{equation}
			u \left( x \right) = \sum_{i = 0}^{6} d_i x^i,
		\end{equation}
		where
		{\allowdisplaybreaks
		\begin{align}
			d_0 &= \gls{delta},	\\
			d_1 &= \frac{e_2}{e_1} P + \frac{e_2}{e_1} c_0 + \frac{e_2}{2 e_1} c_1 + \frac{e_2}{3 e_1} c_2,	\\
			d_2 &= -\frac{e_2}{2 e_1} c_0,	\\
			d_3 &= \frac{P}{3 e_1} + \frac{1}{3 e_1} c_0 + \frac{1 - e_2}{6 e_1} c_1 + \frac{1}{9 e_1} c_2,	\\
			d_4 &= -\frac{1}{4 e_1} c_0 - \frac{e_2}{12 e_1} c_2,	\\
			d_5 &= -\frac{1}{10 e_1} c_1,	\\
			d_6 &= -\frac{1}{18 e_1} c_2.
		\end{align} }% end of \allowdisplaybreaks
	\end{subequations}
\end{case}	\text{}
%
% An exponential q(x)
\begin{case} \label{cas:q_Exponential}
	Consider also \cref{prb:Normalized_Problem}, with an exponential $q \left( x \right)$, i.e.,
	\begin{subequations} \label{eqn:q_Exponential_Problem}
		\begin{align}
			-\frac{d}{dx} \left[ q \left( x \right) \frac{du}{dx} \right] &= f_{n} \left( x \right),
				\quad x \in \left( 0, 1 \right), \label{eqn:q_Exponential_Diff_Eq}	\\
			\gls{q} &= \exp \left( - \alpha x \right),  \label{eqn:Exponential_q_Func} \\
			\gls{fDegreen} &= \sum_{i = 0}^{n}c_{i}x^{i}, \label{eqn:Exponential_Polynomial_Load}
		\end{align}
		with boundary conditions
		\begin{align}
			u \left( x \right) \Bigg \vert_{x = 0} &= \gls{delta},  \label{eqn:q_Exponential_Left_Boundary} \\
			\left[ q \left( x \right) \frac{du}{dx} \right]_{x = 1} &= P. \label{eqn:q_Exponential_Right_Boundary}
		\end{align}
	\end{subequations}
	The implementation requires $\alpha \neq 0$ because the analytical formula below contains inverse powers of $\alpha$; $\alpha=0$ is represented by the constant model.
	Again, we solve this for a polynomial load function up to second degree, i.e,
	\begin{subequations}	\label{eqn:q_Exponential_Solution}
		\begin{equation}
			u \left( x \right) = d_0 + \exp \left( \alpha x \right) \sum_{i = 1}^{4} d_i x^{i-1},
		\end{equation}
		where
		{\allowdisplaybreaks
		\begin{align}
			d_0 &= \gls{delta} - \frac{P}{\alpha} - \left( \frac{1}{\alpha} + \frac{1}{\alpha^2} \right) c_0
					+\left( - \frac{1}{2 \alpha} + \frac{1}{\alpha^3} \right) c_1
						-\left( \frac{1}{3 \alpha} + \frac{2}{\alpha^4} \right) c_2,	\\
			d_1 &= \frac{P}{\alpha} + \left( \frac{1}{\alpha} + \frac{1}{\alpha^2} \right) c_0
							+\left( \frac{1}{2 \alpha} - \frac{1}{\alpha^3} \right) c_1
							+ \left( \frac{1}{3 \alpha} + \frac{2}{\alpha^4} \right) c_2, \\
			d_2 &= -\frac{1}{\alpha} c_0 + \frac{1}{\alpha^2} c_1 - \frac{2}{\alpha^3} c_2, \\
			d_3 &= -\frac{1}{2 \alpha} c_1 + \frac{1}{\alpha^2} c_2, \\
			d_4 &= -\frac{1}{3 \alpha} c_2.
		\end{align}}	% end of \allowdisplaybreaks
	\end{subequations}
\end{case}
%
% A constant q(x)
\begin{case} \label{cas:q_Constant}
	Finally, consider \cref{prb:Normalized_Problem}, with a constant rigidity, i.e.,
	\begin{subequations} \label{eqn:q_Constant_Problem}
		\begin{alignat}{2}
			-\frac{d}{dx} \left[ q \left( x \right) \frac{du}{dx} \right] &= f_{n} \left( x \right),
				\quad && x \in \left( 0, 1 \right),	\label{eqn:q_Constant_Diff_Eq}	\\
			\gls{q} &= e_0,  \label{eqn:Constant_q_Func} \\
			\gls{fDegreen} &= \sum_{i = 0}^{n}c_{i}x^{i}, \quad && n = 0, 1, 2, \label{eqn:Constant_Polynomial_Load}
		\end{alignat}
		with boundary conditions
		\begin{align}
			u \left( x \right) \Bigg \vert_{x = 0} & = \gls{delta},  \label{eqn:q_Constant_Left_Boundary} \\
			\left[ q \left( x \right) \frac{du}{dx} \right]_{x = 1} &= P. \label{eqn:q_Constant_Right_Boundary}
		\end{align}
	\end{subequations}
The constant rigidity coefficient must satisfy $e_0>0$.
The solution is given by
	\begin{subequations}	\label{eqn:q_Constant_Solution}
		\begin{equation}
			u \left( x \right) = \sum_{i = 0}^{4} d_i x^i,
		\end{equation}
		where
		{\allowdisplaybreaks
		\begin{align}
			d_0 &= \gls{delta}, \\
			d_1 &= \left( P + c_0 + \frac{c_1}{2} + \frac{c_2}{3} \right) \frac{1}{e_0}, \\
			d_2 &= -\frac{c_0}{2e_0}, \\
			d_3 &= -\frac{c_1}{6e_0}, \\
			d_4 &= -\frac{c_2}{12e_0}.
		\end{align}}	% end of \allowdisplaybreaks
	\end{subequations}
\end{case}
In \cref{lst:Def_Problem} we have included the MATLAB code for implementation of all analytical solutions discussed here.
%